\scanpage{44}\mspage{43}
Cette topologie est plus fine que $\tau\deleted{'}$, i.e. si $X$ est fermé dans $\tau$, il satisfait ces conditions : clair. Il induit $\tau'_0$, i.e. si $X\subset T$, il satisfait ces cond. sss il est fermé pour $\tau'_0$ : clair. Donc $\tau'\in\underline{\mathscr{T}}$. Prouvons que c'est un plus grand élément (top la moins fine). i.e que pour \deleted{\unsure{tout|tous}} $\tau''\in\underline{\mathscr{T}}$, les fermés pour $\tau'$ sont des fermés pour $\tau''$. En effet, \unsure{on a}
\[
X=B\deleted{\cap}\cup\overline{X\capU}^{(\tau)}
\]
$B$ étant fermé dans $T$. (pour $\tau'_0$\inserted{\deleted{\unsure{donc}} pour $\tau^{\prime\prime\unsure{,}}\mid\overset{\sim}{T}$}), et $T$ étant \uline{fermé} dans $\deleted{\unsure{\tau|T|B}}S$ pour $\tau$ donc \inserted{a fortiori} pour $\tau''$, $B$ est fermé dans $S$ pour $\tau''$.
\margin{(\unsure{C'est|car} la première fois qu'on utilise $T$ fermé dans $S$}
De $=$ $\overline{X\cap U}^{(\tau)}$ étant fermé pour $\tau$, l'est pour $\unsure{?|;|:}\tau''$, donc $X=B\cup\overline{X\cap U}^{(\tau)}$ l'est aussi. Cela prouve a) b) c).

Prouvons d)\inserted{(\unsure{le début|ci dessus} déjà \unsure{prouvé|fini|fait})}, en nous rappelant que pour $x,y\in T$,
\[
\begin{aligned}
x\underset{\tau'}{\leq}y&\Longleftrightarrow\overline{x}^{(\tau')}\subset\overline{y}^{(\tau')}\\
&\Longleftrightarrow x\in\overline{y}^{(\tau')}.
\end{aligned}
\]
Comme $\tau'\mid U\ne\tau\mid U$, et $\tau'\mid T=\tau'_0$, les cas $x,y\in T$ et $x,y\in U$ sont déjà clairs. Reste le cas où soit $x\in U,y\in T$, soit l'inverse. Le premier cas ne peut se produire, car si $y\in T$, $\overline{y}^{(\tau')}\inT$ ($T$ \uline{fermé} dans $S$ pour $\tau'$) donc $x\notin\overline{y}^{(\tau')}$ \deleted{\unsure{puisque|puisq.}} si $x\in U$. Donc reste le cas $[x\in T,]$ $y\in U$, prouver $x\underset{\tau'}{\leq}y\Longleftrightarrow x\underset{\tau'}{\leq}y$
\margin{NB c'est vrai sans supposer $x\inT$}
